\iffalse
This section has the framework

1. Theory of the Molecular Dynamics 
2. Optical Physics and How we get signal (iSCAT)
3. Biological Applications 

Appendix 1.
- Basic E and M
- Charged surface potentials
- Stern Layers, Poisson Boltzmann ( Very short )
- Rayleigh Scattering Derivation (?)

Appendix 2
- Stochastic Processes
- Stationary solutions to the Smoluchowski Equation (Boltzmann Distribution)

Appendix 3.
- Fitting Methods used here
- L-BGFS

Appendix 4.
- Grid Traversal Algorithm
- Grid Search Method (?)
\fi

Biology is mostly nanoparticles. Nanoparticles move around stochastically under the rules of brownian motion. To be able to measure the particle's characteristics, one can relate their motion and stochasticity to their properties. This chapter will discuss the phenomena of Brownian motion and how it relates to particle properties, environmental ion concentration, viscosity, hydrodynamic boundaries, and confinement. To directly measure particle position, one can use optical methods, specifically, optical scattering measurements. In this chapter, the theoretical background regarding the particle's stochastic motion is presented, with particular focus on how the processes evolve in confined geometries. There will be a discussion on relating optical signals to their motion, as well as optical methods for particle detection. 




Methods that are able to detect single nanoparticles are known as single-molecule measurements. In the field of optical microscopy, single-molecule methods are often referred to as single particle tracking (SPT). This chapter aims to direct the reader through the theoretical context of the work done. First, Brownian motion is discussed, then optical preliminaries for SPT, and lastly, biological motivation is touched upon to continue the introductory words.



\section{Brownian Motion}
\iffalse
- Introduce the concept of Brownian Motion
- Histroy of brownian Motion
- applications of simple

- discuss an outlook of what has been done since then
- outline of this section
\fi

The physical phenomenon of diffusion can be rigorously described as the emergent behavior atomic particle undergoing random collisions with members of the surrounding environment. In the microscopic limit, where the radius of a particle is sufficiently small and thermal fluctuations dominate its dynamics, this erratic motion is referred to as Brownian motion \cite{brown_xxvii_1828}. 

\subsection{Classical Description of Einstein's Brownian Motion}

The classical description of Brownian motion is concerned with how the molecuar-kinetic theory of heat reates with microscopically visible size suspensions in liquid perform movements such that they can be viewed by a microscope. 


\subsection{Langevin Dynamics}

- Brownian motion can alternatively be described using a Langevin Model. 
- introduce langevin equation
- velocity correlation
- Arrive at the Fluc-Dissipation
- relate it to Einstein's formulation and show that given a stokes drag (sphere)

\begin{equation}
    D = \frac{k_bT}{\gamma_0}
\end{equation}

\subsection{Smoluchowski Equation}

- Since Langevin is kind of convient, how do we acctualy link the two methods?

- Fokker plank introdution motivated from the master equation
- - ideally summarize van kampen. 

- arrive at the general Smoluchowski equation for a random walker
- simple plots for descriptive purposes.

\section{Confined Diffusion Processes}

Although diffusion is interesting, there is a wealth of complexity in confinement due to Hydrodynamic effects
- quote 111 years here

- outline of this seciton

\subsection{Hydrodynamic Confinement}
- confinemnt from Hydrodynamics
- frictional slowing
- ...

\subsection{Geometric Confinement}
- i need to check references on if this is relate

\subsection{Electrostatic Potentials}

\subsection{Example: Polystyrene Bead}


\section{Interferometric Scattering Principles}

This section is devoted to the description of physical prinicples associated with light scattering techniques. These techniques are widely used in microscopy, as I have used them in measuring where my parteicles are. Scattering measurements are of interest as they allow for label-free imaging with a high signal to noise and high temporal resolution. 

We begin with the basics of relating light to electromagnetsm, starting with the Maxwell and Helmholtz relations. We continue discussing the scattering of a various size range of spheres, which is the approximation we take in this thesis. We then discuss the various abberations and Considerations before introducting iSCAT microscopy. We tease the experimental apparatus used in tthis work and compare the method to other methods currently used in the field. 

\subsection{Maxwell's Equations and field propagations equations}

The photon is the fundamental quanta of the electric field as stated in QFT (cite). It has been shown that because of this, the field equations of light obey the fundamental principle that states
\begin{equation}
    I = |\vec E(\vec x, t)|^2
\end{equation}.
Thus, it is paramount that we understand how the propagation of the electric fields are described. Luckily for us, this has already been solved in the year 18XX (cite). We present the Maxwell equations in their macroscopic form, which include terms from the material constants.
\begin{definition}
    Maxwell's Equations in macroscopic form. Let $\vec x$ be the position vecotr, $t$ the time scalar. The Maxwell equations are defined as the following set of relations; (CITE ZANGWILL, Scattering small particles, Jackson if needed).
    \begin{align}
        \nabla \cdot \vec D(\vec{x},t) &= \rho_f(\vec{x},t) \\
        \nabla \times \vec E(\vec{x},t) + \frac{\partial \vec B(\vec{x},t)}{\partial t} &= 0 \\
        \nabla \cdot \vec B(\vec{x},t) &= 0 \\
        \nabla \times \vec H(\vec{x},t) - \frac{\partial \vec D(\vec{x},t)}{\partial t} &= \vec J_f(\vec{x},t)
    \end{align}
    In addition, for ease of access, the definitions for the Auxillary fields, and scalar potentials are written. We also add the guage invariance;
    \begin{align}
        \vec D(\vec{x}, t) &= \varepsilon_0 \vec E(\vec{x},t) + \vec P (\vec{x},t) \\
\vec H (\vec{x},t) &= \frac{1}{\mu_0} \vec B(\vec{x},t) - \vec M(\vec{x},t)
    \end{align}
    To fully describe the material that the field lies in, we also require constitutive relations;
    \begin{align}
        \vec J_f(\vec{x},t) &= \mathbf{\sigma} \vec E (\vec{x},t)\\
        \vec B(\vec{x},t) &= \mathbf{\mu} \vec H (\vec{x},t)\\
        \vec P(\vec{x},t) &= \varepsilon_0\mathbf{\chi} \vec E (\vec{x},t)\\
    \end{align}
    Where the boldface is reserved for tensor quantities. 
\end{definition}

- List shit that we need to know before calcualting the Scattering from a sphere.

- Helmholtz Equation
- Fresnel Equations
- Polarization
- General Scattering Transfer Matrices
\subsection{Electromagnetic Scattering of of Homogenous Spheres}

To As mentioned, we are tasked with getting the light from the scattered molecules. We approximate them as spheres. Thus, we need to solve the solutions of scattered light from spheres.

We start by using the wave equation for the Electric and Magnetic Fields
\begin{equation}
    \nabla^2 \vec E + k^2 \vec E = 0, \quad \vec H = \frac{1}{i \omega \mu_0} \nabla \times \vec E
\end{equation}

The solutions of these quantities will be harmonics of the laplacian in the coordinate basis that we choose. For isotropic scattering, the spherical basis is the one of choice. The solutions to this are known as the Vector Spherical Harmonics. Let $\psi$ be a solution to the scalar wave equation
\begin{equation}
    \nabla^2 \psi + k^2 \psi = 0
\end{equation}, then we can form a generator of the general solution by taking the simplest form of solutions and summing all components. Citing Jackson, we find that the solution to this form is given by the seperable form;
\begin{equation}
    \begin{aligned}
        \psi(r, \theta, \phi) &= \Phi(\phi) \Theta(\theta) R(r) \\
        \Phi(\phi) & \in \{\cos(m\phi), \sin(m\phi)\} \\
        \Theta(\theta) &= P^m_n(\cos\theta)\\
        R(r) &\in \{j_n, y_n, j_n + i y_n, j_n - i y_n\} \\
    \end{aligned}
\end{equation}
Where $j_n, y_n$ are the spherical Bessel functions, and their complex additions are the hankel functions. Thus, the vector spherical harmonics are best described by seperating these solutions into even and odd solutions as their prefactor will be $\cos$ or $\sin$ respectively. 
\begin{equation}
    \begin{aligned}
        \vec{M}_{emn} = \vec \nabla \times (\vec r \psi_{emn}), \quad \vec{M}_{omn} = \vec \nabla \times (\vec r \psi_{omn}) \\
        \vec{N}_{emn} = \frac{1}{k}\vec \nabla \times \vec{M}_{emn}, \quad \vec{N}_{omn} = \frac{1}{k}\vec \nabla \times \vec{M}_{omn} \\
    \end{aligned}
\end{equation}
These are found in Eqns 4.17 to 4.20 in Bohren. 

The incident light can be assumed to be polarized along the x axis without loss of generality. The general form of these equation will then be the sum of each of the spherical vector harmonics (SVH). subject to the energy conservation boundary condition
\begin{equation}
    (\vec E_i + \vec E_s - \vec E_1) \times \hat r = (\vec H_i + \vec H_s - \vec H_1) \times \hat r = 0
\end{equation}

Citing Bohren, and stating that using orthongality of the SVHs, we have that the exact solution, which is known as the Mie Theory form of scattering off of a sphere, is given below;

\begin{equation}\label{eqn:lorenzmie}
    \begin{aligned}
     \vec E_i &= E_0 e^{ikr\cos \theta}\hat x \\    
     \vec E_s &= \sum_{n=1}^\infty E_n(i a_n \vec{N^{(1)}}_{e1n} - b_n \vec{M^{(1)}}_{o1n})\\
     \vec E_1 &= \sum_{n=1}^\infty E_n(i c_n \vec{M^{(1)}}_{o1n} - d_n \vec{N^{(1)}}_{e1n})     
    \end{aligned}
\end{equation}

Since we are only interested in the scattered field, the constants that we care about are written below;
\begin{equation}
    \begin{aligned}
        a_n &= \frac{m\Psi(mx)\Psi'(x) - \Psi(x)\Psi'(mx)}{m\Psi(mx)\xi'(x) -\xi(x) \Psi'(mx)}\\
        b_n &= \frac{\Psi(mx)\Psi'(x) - m\Psi(x)\Psi'(mx)}{\Psi(mx)\xi'(x) - m \xi(x) \Psi'(mx)}\\
    \end{aligned}
\end{equation}

where $\Psi, \xi$ are the Riccatti-Bessel functions. 

In the limit of a small particle from far away, we can mae some cancellations in these constants to simplify the work. Notably, the hankel function approches the limit of a spherical wave,
\begin{equation}
    h_n(kr) \xrightarrow{kr \gg 1} (-i)^{n+1}\frac{e^{ikr}}{kr}
\end{equation}

Im this regime, we have that $a_1 \approx \frac{-2i}{3}(ka)^3 \frac{m^2 -1 }{m^2 + 2}$, and $b_1 \approx 0$. Thus for $n=1$ we have that $\pi_1 = 1, \tau_1 = \cos \theta$, thus
\begin{equation}
    \begin{aligned}
        M^{(3)} &\approx (-i)^{n+1}\frac{e^{ikr}}{kr} \bracket{\cos\phi \pi_n \hat \theta - \sin\phi \tau_n \hat \phi} \\
        N^{(3)} &\approx (-i)^{n+1}\frac{e^{ikr}}{kr} \bracket{\cos\phi \pi_n \hat \theta - \sin\phi \tau_n \hat \phi} \\
    \end{aligned}
\end{equation}

can be written in the Rayleigh approximation simply and in the form of the full scattering field;

\begin{equation}
    \vec E_s = \frac{i E_0 e^{ikr}}{kr}\sbracket{\hat \theta \cos \phi \sum_{n=1}^\infty \frac{2n +1}{n(n+1)}(a_n \tau_n + b_n \pi_n) + \hat \phi \sin \phi \sum_{n=1}^\infty \frac{2n +1}{n(n+1)}(a_n \pi_n + b_n \tau_n)}
\end{equation}

Spoiling the suprise, after plugging in the terms for $n=1$ we have the following simplification;

\begin{equation}
    \begin{aligned}
            \vec E_{s, \perp} &=  \vec E_{i, \perp} \frac{k^2 \alpha e^{ikr}}{r}\\
            \vec E_{s, \parallel} &=  \vec E_{i, \parallel} \frac{k^2 \alpha e^{ikr}}{r} \cos \theta\\
    \end{aligned}
\end{equation}

Then, to find the scattering intensity field, and noting $k = \frac{2\pi}{\lambda/n}$, we have the final scattering field in the Rayleigh regime;
\begin{equation}
    \vec E_s = \vec E_0 e^{ikr}\frac{\sqrt{8}\alpha \pi^2 n^2_m}{\lambda^2 r} \sqrt{1 + \cos \theta}
\end{equation}

- Draw some of the iPSFs for Mie theory, from all angles.

- Set the stage, what do we need to solve, how are we going to do it
- Solutions to the wave equation in spherical coordinates 
- Mie Theory 
-- Examples of Mie spheres
-- What size range ?
- RDG Theory
-- What size range ?

- Small comparison from far away, medium (mie), and small (RDG), atomic (too small)

\subsection{Types of Microscopy}

Dark field - microscopy where the background is not used
widefield - microscopy where the full field of view is used
holography - microscopy with a reference field (?)

\subsection{Scattering Microscopy}

As shown, the description of the scattered light from a sphere is full described by Eqn. \ref{eqn:lorenzmie}, known as the Lorenz-Mie scattering theory (CITE). In the limit of particles on the order of the size of the wavelength of light, the Rayleigh model is an easier, equivalent method derived from Mie theory. 

To collect the light, there have been three main schemes used in literature; forward, side, and backscattering micrpscopy schemes.

As shown in Fig. (reference PSF), the scattering distribution taken at different projections widely varies. In addition, the amount of light that gets collected in each of these regimes also varies.

\textbf{do i go into the other models?}

In the backscattering scheme, the light collected is the sum of the incident light and the scattered light.

\begin{equation}
    I(x,y) = |\vec E_i + \vec E_s|^2 
\end{equation}

given that the incident wave is known as a polarized plane wave, this equation can be rewritten for the scatterd light;
\begin{equation}
    \begin{aligned}
        I = E_iE_i^* + 2\Re{\left[E_iE_s\right]} + E_sE_s^* \\
    \end{aligned}
\end{equation}
The scattering term will be must lower than that of the 



\dots

- Olsén's thesis



\subsection{Interferometric Light Scattering (iSCAT)}
- Experimental Setup

- manoharan theory and reparam of the main equation
- schematic
- ...

\subsection{Abberations and Image formation}

discuss how we go from reality (focus plane) to image (image plane).

- sphereical abberations
- guoy phase
\dots 

- final kind of check 

theortical 100 nm particle -> scattering field -> abberations 

\section{In-Vivo Considerations}

- Now that we more or less know how BM works, why is it important in cells?

\subsection{why does this matter for biological applications}

\subsection{Examples of Brownian Motion in Biology}

\section{Brownian motion as a method for measurements of size in biology}

\subsection{survey of methods using in this field}
- Optical Methods 
- Bulk: DLS, 
- Single Molecule: SPT, NTA, FCS, 

\subsection{how confinemen effects these measurements}

- Literature on why these methods are lacking

\subsection{Probing of surface interactions with these methods}






